\section{A Role-Sensitive Framework}
\label{sec:framework}

The framework has three parts: the \emph{role} a model plays, the \emph{level}
at which evidence about it is gathered, and the \emph{claim} that the evidence
is meant to support. A fourth consideration, the \emph{tier} of risk, scales
how much evidence is proportionate. Together they answer the question that a
single standard cannot: what would count as adequate evaluation for
\emph{this} use of \emph{this} model in \emph{this} study.

\subsection{Nine Roles}
\label{sec:roles}

Figure~\ref{fig:roles} and Table~\ref{tab:roles} set out nine functional
roles. Six are roles in the robot's interaction loop. Three are roles in the
research pipeline, where a model is used to study the interaction and is never
encountered by a participant. The distinction matters because the two groups
fail differently. An error in the loop reaches a person. An error in the
pipeline reaches the literature.

% Figure 1: functional roles of generative models in an HRI pipeline
\begin{figure}[t]
\centering
\resizebox{\textwidth}{!}{%
\begin{tikzpicture}[
    font=\sffamily\footnotesize,
    role/.style={draw=black!70, rounded corners=2pt, minimum height=11mm,
                 text width=23mm, align=center, fill=black!4, inner sep=2pt},
    inst/.style={draw=black!70, dashed, rounded corners=2pt, minimum height=9mm,
                 text width=27mm, align=center, fill=white, inner sep=2pt},
    actor/.style={draw=black!70, thick, minimum height=11mm, text width=15mm,
                  align=center, fill=white},
    flow/.style={-{Latex[length=2mm]}, semithick, black!75},
    obs/.style={-{Latex[length=2mm]}, semithick, dashed, black!60},
    band/.style={font=\sffamily\scriptsize\itshape, text=black!70}]

% --- actors --------------------------------------------------------------
\node[actor] (human) at (0,0) {Person(s)};
\node[actor] (human2) at (14.6,0) {Person(s)};

% --- in-robot roles ------------------------------------------------------
\node[role] (perc) at (2.65,0) {\textbf{R1 Perceiver}\\interprets people\\and scenes};
\node[role] (rep)  at (5.45,0) {\textbf{R2 Representer}\\embeds, retrieves,\\remembers};
\node[role] (plan) at (8.25,0) {\textbf{R3 Decider}\\plans, selects,\\times, defers};
\node[role] (lang) at (11.4,1.45) {\textbf{R4 Speaker}\\user-facing\\language};
\node[role] (beh)  at (11.4,0) {\textbf{R5 Behaviour author}\\expressive behaviour,\\code, parameters};
\node[role] (act)  at (11.4,-1.45) {\textbf{R6 Actor}\\end-to-end\\physical action};

\draw[flow] (human) -- (perc);
\draw[flow] (perc) -- (rep);
\draw[flow] (rep) -- (plan);
\draw[flow] (plan.east) -- ++(0.35,0) |- (lang.west);
\draw[flow] (plan.east) -- (beh.west);
\draw[flow] (plan.east) -- ++(0.35,0) |- (act.west);
\draw[flow] (lang.east) -- ++(0.45,0) |- ([yshift=2.5mm]human2.west);
\draw[flow] (beh.east) -- (human2.west);
\draw[flow] (act.east) -- ++(0.45,0) |- ([yshift=-2.5mm]human2.west);

% robot boundary
\begin{scope}[on background layer]
  \node[draw=black!55, thick, rounded corners=4pt, fill=black!1,
        fit=(perc)(rep)(plan)(lang)(beh)(act), inner sep=3.2mm] (robot) {};
\end{scope}
\node[band, anchor=south west] at ([xshift=1mm,yshift=0.2mm]robot.north west)
     {Robot architecture (models deployed in the interaction loop)};

% end-to-end bypass (VLA)
\draw[flow, black!45] (perc.south) |- (act.west |- 0,-2.05) -| ([xshift=-6mm]act.south);
\node[band] at (6.1,-1.86) {end-to-end vision--language--action policies collapse R1--R3 into R6};

% --- research-instrument roles ------------------------------------------
\node[inst] (anno) at (3.3,-3.75) {\textbf{R7 Annotator}\\codes transcripts,\\video, logs};
\node[inst] (judge) at (7.3,-3.75) {\textbf{R8 Judge}\\scores outputs of\\R1--R6};
\node[inst] (sim) at (11.3,-3.75) {\textbf{R9 Simulator}\\synthetic users,\\partners, raters};

\begin{scope}[on background layer]
  \node[draw=black!55, thick, dashed, rounded corners=4pt,
        fit=(anno)(judge)(sim), inner sep=3.2mm] (lab) {};
\end{scope}
\node[band, anchor=north west] at ([xshift=1mm,yshift=-0.2mm]lab.south west)
     {Research pipeline (models used to study the interaction)};

\draw[obs] (anno.north) -- (anno.north |- robot.south);
\draw[obs] (judge.north) -- (judge.north |- robot.south);
\draw[obs] (sim.north) -- (sim.north |- robot.south);
\end{tikzpicture}}
\caption{Nine functional roles that language, vision, and action models play in
HRI research. Solid boxes (R1--R6) are roles inside the robot's interaction
loop; dashed boxes (R7--R9) are roles inside the research pipeline that
observes, scores, or stands in for the interaction. A single model usually
occupies several roles at once, and each role inherits a different evaluation
tradition (Table~\ref{tab:roles}).}
\Description{A block diagram. A person on the left feeds into a chain of three
boxes inside a robot boundary labelled R1 Perceiver, R2 Representer and R3
Decider. The Decider fans out to three output boxes, R4 Speaker, R5 Behaviour
author and R6 Actor, each of which points to a person on the right. A grey
arrow bypasses the chain from the Perceiver to the Actor to indicate
end-to-end vision-language-action policies. Below the robot boundary, a dashed
boundary labelled research pipeline contains three dashed boxes, R7 Annotator,
R8 Judge and R9 Simulator, with dashed arrows pointing up at the robot.}
\label{fig:roles}
\end{figure}


\begin{table}[t]
\caption{Functional roles of generative models in HRI, the evaluation
tradition each inherits, and the component-level evidence that the tradition
treats as the minimum.}
\label{tab:roles}
\small
\begin{tabularx}{\textwidth}{@{}L{2.25cm}L{3.6cm}L{2.9cm}Y@{}}
\toprule
Role & Typical use in HRI & Inherited tradition & Minimum component-level evidence \\
\midrule
\role{1} Perceiver & Infers affect, engagement, intent, or wellbeing from
speech, text, or video; describes scenes & Classification; clinical and
behavioural assessment & Per-class performance against multiply rated human
reference; agreement relative to the human ceiling; subgroup error rates \\
\role{2} Representer & Embeds utterances or images for retrieval, memory,
similarity, clustering, user modelling & Information retrieval; construct
validity & Retrieval accuracy; stability across embedding models and seeds;
null baselines; human check of derived structure \\
\role{3} Decider & Selects plans, skills, or tools; decides when to speak,
help, interrupt, ask, or defer & Planning; judgement and decision-making &
Agreement with expert decisions; constraint violations; calibration of
abstention and help-seeking \\
\role{4} Speaker & Generates language delivered to a person as speech or text
& Human evaluation of generated text and dialogue & Rubric-based ratings by
trained raters with reliability; error taxonomy with severity; audit of
delivered outputs \\
\role{5} Behaviour author & Generates gestures, facial expressions, motion
parameters, or robot programs & Behaviour-perception studies; program
verification & Recognition by naive observers against chance; executability
and safety checks; comparison with expert-authored behaviour \\
\role{6} Actor & Maps observation and instruction directly to motor commands &
Trial-based policy evaluation; safety engineering & Success rates with
intervals over specified conditions; failure taxonomy; safety-envelope
violations; reliably adjudicated outcomes \\
\midrule
\role{7} Annotator & Codes transcripts, video, or logs produced by a study &
Content analysis & Agreement with independent human coders on a held-out
subset, alongside human--human agreement \\
\role{8} Judge & Scores outputs of \role{1}--\role{6} in place of human raters
& Meta-evaluation & Judge--human agreement per construct; bias probes;
independence from the system judged \\
\role{9} Simulator & Stands in for users, partners, or raters & Simulation
validity & Fidelity against human data on the properties relied upon;
restricted to piloting and stress-testing \\
\bottomrule
\end{tabularx}
\end{table}

Three properties of the taxonomy are worth stating. First, roles are
functional, not architectural. A single call to a multimodal model can occupy
several roles at once, as when a model receives a camera frame and a
transcript, decides that the user is distressed (\role{1}), chooses to offer
support instead of continuing the task (\role{3}), and writes the utterance
(\role{4}). The evaluation obligations of all three roles apply, even though
there is one prompt and one output. End-to-end vision-language-action
policies make the same point more sharply: they collapse perception,
representation, and decision into action, which removes the intermediate
outputs that the first three traditions evaluate and leaves only behaviour to
be judged.

Second, the taxonomy makes concrete the contrast with which this article
began. A model in the Speaker role produces open-ended text that a person
hears, so the relevant standard is human evaluation of generated language:
is it relevant, correct, safe, and transparent? A model in the Representer or
Perceiver role produces a vector or a label that a person never sees, so the
relevant standard is that of a measurement instrument or a classifier: is it
accurate against a reliable reference, stable, and unbiased across groups?
Applying the first standard to the second role produces vague reassurance
(``the outputs looked reasonable''); applying the second to the first produces
false precision (an accuracy figure for something that has no single correct
answer).

Third, the research-pipeline roles are included because they are where
generative models have entered HRI most quietly. A study with a fully
scripted robot can still depend on a language model that coded the
transcripts, a vision-language model that scored engagement from video, or
simulated users that were used to tune the protocol. The robot in such a
study is not ``LLM-powered'', and the paper may not appear in any review of
LLM-based HRI, yet its findings are conditional on the model's validity as an
instrument.

\subsection{Four Levels of Evidence}
\label{sec:levels}

Evidence about a model in a role can be gathered at four levels
(Table~\ref{tab:levels}).

\begin{description}[leftmargin=1.2em, style=sameline]
\item[L1 Component.] The model's outputs are evaluated against the standard of
the role's tradition, offline or from logs: are the labels accurate, the
utterances appropriate, the plans executable, the trajectories successful?
\item[L2 System.] The model is evaluated as integrated in the robot and
running in real time: latency, propagation of speech-recognition and
perception errors, consistency between what the model says and what the body
can do, behaviour under network loss, recovery from failure.
\item[L3 Interaction.] People interact with the robot and outcomes are
measured in them: perceptions, behaviour, task performance, learning,
wellbeing. This is the level of the HRI user study.
\item[L4 Deployment.] The robot is used over time outside the laboratory:
drift in model behaviour, novelty effects, accumulation of memory and
personalisation, rare harms, and the organisational context of use.
\end{description}

\begin{table}[t]
\caption{Questions answered at each level of evidence, with typical measures.
The levels are complementary: evidence at one level does not substitute for
evidence at another.}
\label{tab:levels}
\small
\begin{tabularx}{\textwidth}{@{}L{1.9cm}L{4.3cm}Y@{}}
\toprule
Level & Question & Typical measures \\
\midrule
L1 Component & Are the model's outputs adequate for the role? & Agreement
with human reference; rubric ratings; accuracy, F1, calibration; success rate
over trials; error and hallucination rates; run-to-run and prompt-to-prompt
consistency \\
L2 System & Does the integrated robot behave as specified in real time? &
End-to-end response latency distribution; word error rate and its downstream
effect; capability-consistency violations; fallback and timeout rates;
intervention rate of a supervising human \\
L3 Interaction & What happens to the people who interact with it? &
Validated self-report scales; behavioural measures; task outcomes;
manipulation checks; qualitative accounts; adverse events \\
L4 Deployment & Does it keep working, and for whom, over time? & Longitudinal
outcome trajectories; usage and attrition; drift under model updates; incident
reports; subgroup differences in benefit and harm \\
\bottomrule
\end{tabularx}
\end{table}

The central principle is that \emph{evidence does not transfer between levels
without an argument}. Two inferential errors follow from ignoring it, and both
are common.

The \emph{upward leap} takes component-level evidence, often not even the
authors' own, as support for interaction-level claims. A model that leads a
public benchmark, or that the authors found to give sensible answers when
prompted at a desk, is assumed to be adequate when embodied in a robot,
addressed in noisy speech by an older adult, and constrained to answer within
a conversational pause. Benchmarks rarely sample that context, and the
system-level properties that decide whether an interaction succeeds, such as
latency and error propagation, are not properties of the model at all.

The \emph{downward leap} takes interaction-level evidence as support for
component-level adequacy. Participants rated the robot as empathic,
trustworthy, and useful, so its utterances are taken to have been accurate
and safe. This inference fails for a specific reason: participants are poorly
placed to detect the errors that matter most. A fluent, confident, incorrect
answer is by construction one that the listener does not recognise as
incorrect. Favourable questionnaire scores are compatible with a robot that
fabricated facts, overstated its capabilities, or gave advice a professional
would not have given. Only an audit of the outputs that participants actually
received can establish what they received.

\subsection{Five Kinds of Claim}
\label{sec:claims}

Evaluation serves a claim, and the same data can be adequate for one claim
and inadequate for another. We distinguish five.

\begin{description}[leftmargin=1.2em, style=sameline]
\item[Capability.] The model can perform the role to a stated standard under
stated conditions (``the model classified engagement with $F_1=.78$ against
two trained coders'').
\item[Comparison.] One configuration performs differently from another on a
stated outcome (``participants disclosed more to the robot with the adaptive
prompt than to the robot with the static prompt'').
\item[Substitution.] The model can replace a human, or a more expensive
procedure, for a stated purpose (``the model can replace the second coder'';
``the autonomous robot can replace the wizard'').
\item[Safety.] The frequency and severity of harmful outputs or actions are
below a stated bound under stated conditions.
\item[Generalisation.] A finding obtained with one model, prompt, population,
or setting holds for others.
\end{description}

Substitution and safety claims carry the heaviest burden and are the most
often under-supported. A substitution claim is a claim of equivalence or
non-inferiority, and a non-significant difference between model and human in a
small sample is not evidence of either (Section~\ref{sec:equivalence}). A
safety claim concerns the tail of a distribution, and a mean rating says
little about tails. Generalisation claims about ``LLMs'' or ``VLAs'' as a
class, made on the basis of one model at one point in time, are best avoided
or explicitly hedged; what a study shows is that one configuration behaved in
one way.

\subsection{Three Tiers of Proportionate Evidence}
\label{sec:tiers}

Not every use of a model deserves the same scrutiny, and a standard that
demanded maximal evidence everywhere would be ignored everywhere. We propose
that the evidence expected of a study scale with three questions: who is
\emph{exposed} to the model's output and how directly, what is the worst
plausible \emph{consequence} of an inadequate output and can it be reversed,
and how much \emph{oversight} stands between the output and the person.
Table~\ref{tab:tiers} gives three tiers. A study is assigned the highest tier
that any of its roles reaches.

\begin{table}[t]
\caption{Three tiers that scale the expected evidence to exposure,
consequence, and oversight.}
\label{tab:tiers}
\small
\begin{tabularx}{\textwidth}{@{}L{1.8cm}L{4.6cm}Y@{}}
\toprule
Tier & Defining conditions & Expected in addition to the tier below \\
\midrule
I Exploratory & Outputs are reviewed by a researcher before use, or reach
only informed adult participants in low-stakes tasks under supervision;
errors are harmless and reversible; claims are of capability or feasibility &
Role, level, and claim declared; model and configuration fully reported;
outputs logged and a sample inspected and described; limits on the claim
stated \\
II Standard & Outputs reach participants unreviewed in a controlled study, or
model-produced labels and scores feed the study's inferential statistics;
claims are comparative & Component-level validation against a human reference
with reliability statistics; repeated-run consistency; audit of delivered
outputs; validated automated evaluators; uncertainty intervals \\
III High-stakes & Outputs reach children, patients, older adults, or people in
distress; or drive physical action near people; or operate unsupervised over
time; or support substitution or safety claims & Domain-expert raters;
pre-registered evaluation plan; subgroup and failure-tail analysis; adversarial
and stress testing; system-level safety measures with logged human
intervention; monitoring plan for deployment \\
\bottomrule
\end{tabularx}
\end{table}

The tiers are deliberately coarse. Their purpose is to give authors and
reviewers a shared starting point for the question ``is this enough?'', and to
make explicit that the answer depends on the role and the stakes and not on
which model was used. A late-breaking report that explores whether a
vision-language model can describe a robot's workspace need not meet the bar
set for an autonomous robot that conducts wellbeing conversations with
adolescents; but the second study should not be reviewed as if it were the
first.
